% ------------------------------------------------------------------------------
% Kata Pengantar
% ------------------------------------------------------------------------------
\begin{center}
\textbf{\large KATA PENGANTAR}\\
\end{center}
\vspace{0.3cm}

Puji dan syukur penulis panjatkan ke hadirat Tuhan Yang Maha Esa karena atas rahmat dan karunia-Nya, Laporan Tugas Akhir yang berjudul ``\Title'' dapat diselesaikan. Laporan ini disusun untuk memenuhi salah satu syarat memperoleh gelar sarjana pada Program Studi Sarjana \Prodi, Direktorat Kampus Purwokerto, Universitas Telkom. Penulisan Tugas Akhir ini bertujuan untuk menganalisis konsistensi keputusan \textit{LLM-as-a-Judge} terhadap variasi \textit{prompt} dan urutan posisi jawaban dalam Bahasa Indonesia.

Penyusunan Tugas Akhir ini tidak terlepas dari bimbingan, bantuan, dan dukungan berbagai pihak. Oleh karena itu, penulis mengucapkan terima kasih kepada:

\begin{enumerate}\setlength{\itemsep}{2pt}\setlength{\parskip}{0pt}
    \item Bapak Prof. Dr. Suyanto, S.T., M.Sc., selaku Rektor Universitas Telkom.
    \item Bapak Prof. Dr. Z.K. Abdurahman Baizal, S.Si., M.Kom., selaku Direktur Universitas Telkom Direktorat Kampus Purwokerto.
    \item Bapak Aditya Dwi Putro W., S.Kom., M.Kom., selaku Ketua Program Studi Sarjana \Prodi\ Universitas Telkom Direktorat Kampus Purwokerto.
    \item Bapak Andi Prademon Yunus, S.T., M.Sc.Eng., Ph.D., selaku Dosen Pembimbing I yang telah memberikan bimbingan dan arahan selama penyusunan Tugas Akhir ini.
    \item Ibu Nisrina Hanifa Setiono, S.Kom., M.Cs., selaku Dosen Pembimbing II yang telah memberikan bimbingan dan masukan selama penyusunan Tugas Akhir ini.
    \item Kedua orang tua dan keluarga, Bapak dan Ibu dosen, rekan-rekan mahasiswa, serta semua pihak yang telah memberikan doa, bantuan, dan dukungan selama penyusunan Tugas Akhir ini.
\end{enumerate}

Penulis menyadari bahwa laporan ini masih memiliki kekurangan, sehingga kritik dan saran yang membangun tetap penulis harapkan. Semoga Laporan Tugas Akhir ini dapat bermanfaat bagi pembaca.

\vspace{0.5cm}

\noindent\hfill
\begin{tabular}{@{}c@{}}
Purwokerto, \Tanggal \ \Bulan \ \Date\\[1.5cm]
% \includegraphics[height=2cm]{Aset/ttd_mahasiswa.png} % Uncomment untuk TTD
\Author
\end{tabular}
